% Standalone section source: no preamble or external macros are required.
\section{Reverse-suffix response compilation with singular interiors}
\label{sec:streaming-response}

The preceding fixed-suffix response construction leaves an avoidable source
of repeated work: after the scan advances, the next suffix is usually rebuilt
and its entire interior is eliminated again.  We give a different construction
which prepares \emph{all} suffix responses in reverse order.  Its quantitative
gain is a theorem about response compilation on diagrams of bounded frontier.
It supplies neither a universal frontier bound nor a bound on the number of
Khovanov components.  The new implementation is an experimental module; no
recognizer driver dispatches to it automatically.

This construction belongs to the established study of elimination on
graphs of small width.  In particular, F\"urer, Hoppen and Trevisan
give an $O(k^2(m+n))$ rank and determinant algorithm over arbitrary
fields when a suitable width-$k$ decomposition is supplied
\cite{furer-hoppen-trevisan}.  We do not claim a new general
small-width Gaussian-elimination bound.  The contribution here is the
all-suffix determinant response interface, its exact compatibility with
cut-face identifications, and a singular-safe implementation whose
updates can be reused across actual stages of the unknot scanner.

\subsection{The admissible graph filtration}

Fix a checkerboard coloring of the source diagram once for the entire
filtration. In this section, $\alpha$ pairs the two darts of each diagram
edge, and $\rho$ rotates the four darts of each crossing counterclockwise.
For a suffix, its common Tait vertices are black cut-face
fragments, and its edges are the signed crossing edges.  Build suffixes in
reverse: start with the last crossing, then insert the preceding crossing,
and continue.  A crossing initially contributes four new face fragments,
of which exactly two are black.  Whenever a new edge joins an already
present dart, the face permutation adds the two connections
\[
 d\longleftrightarrow\rho(\alpha(d)),\qquad
 \alpha(d)\longleftrightarrow\rho(d).
\]
These connections merge fragments.  They never split a fragment.

A closed fragment is already a complete cycle of the partial face
permutation.  All of its incoming and outgoing permutation arrows are
present; no subsequent insertion can touch it.  Consequently every later
identification occurs among open fragments or fragments just introduced.
Every later signed edge has endpoints among those same fragments.  This
is the essential admissibility condition for retaining an eliminated
interior.  It would be false for arbitrary changes of a diagram or for
forward deletion of an already eliminated crossing.

Keep one distinguished black fragment as a grounding anchor, even if it
closes.  When it merges with another open fragment, follow the merged
fragment.  This leaves at least one retained coordinate at every stage.
If the suffix has $w$ boundary darts, at most $w$ black fragments are open,
so the number $b$ of retained black coordinates obeys
\[
                         1\le b\le w+1.
\]
The extra coordinate is unnecessary when the anchor is already open.
The anchor represents an actual graph vertex; it is not an isolated
vertex added to the graph.

\subsection{A response preserved under arbitrary later attachments}

Work over a field $\mathbb F$.  Starting from a symmetric signed
Laplacian, eliminate only interior coordinates by nonsingular symmetric
principal pivots.  A nonzero diagonal gives a one-coordinate pivot.
If all remaining interior diagonals are zero but an off-diagonal entry
$c$ is nonzero, the principal block
\[
             \begin{pmatrix}0&c\\c&0\end{pmatrix}
\]
is invertible and gives a two-coordinate pivot.  If neither kind exists,
the residual interior block is identically zero.  These alternatives
hold in characteristic two as well as in odd characteristic.

The resulting response consists of a nonzero scalar $\delta$ and
\[
 M=\begin{pmatrix}0_r&R\\R^{\mathsf T}&S\end{pmatrix},
 \qquad R\in\mathbb F^{r\times b},
 \quad S=S^{\mathsf T}\in\mathbb F^{b\times b}.
\]
Here $\delta$ is the product of all eliminated principal determinants.
For any additional graph attached only at the retained coordinates,
and for any identifications among those coordinates, a grounded cofactor
of the augmented original graph equals $\delta$ times the corresponding
grounded determinant obtained from $M$.  To see this, add the exterior
matrix in its own coordinates before eliminating.  Every eliminated
principal block is unchanged, its coupling to new exterior coordinates
is zero, and block Gaussian elimination therefore gives the same Schur
complement and determinant factor.  Terminal identification is the
congruence $M\mapsto T^{\mathsf T}MT$ with $T$ equal to the identity on
interior coordinates.  This congruence commutes with the same interior
elimination.  Grounding a retained terminal also commutes with it.

No invertibility assumption is imposed on the full interior block.
Discarding $0_r$ and $R$ would be unsound.  For example, over
$\mathbb F_2$ the path $0$--$1$--$2$, with terminal endpoints, has the
single interior diagonal $2=0$, but its tree cofactor is one.  Retaining
its one-dimensional radical permits a later two-coordinate pivot when
the path is extended; treating the zero interior pivot as a zero answer
would lose a valid response.

\subsection{The absorbing-zero lemma}

\begin{lemma}[Absorbing zero]
\label{lem:response-zero}
Suppose a residual response has $r\ge b$.  Every grounded cofactor
obtained by any future admissible graph attachments and terminal
identifications is zero over $\mathbb F$.
\end{lemma}

\begin{proof}
Schur complementation of a zero-row-sum matrix still has zero row sums.
Indeed, if
$L=\left(\begin{smallmatrix}P&B\\B^{\mathsf T}&C\end{smallmatrix}\right)$
annihilates the all-ones vector, then
$P\mathbf1+B\mathbf1=0$ and
$B^{\mathsf T}\mathbf1+C\mathbf1=0$, whence
$(C-B^{\mathsf T}P^{-1}B)\mathbf1=0$.
The same is true after terminal congruences and addition of signed
Laplacian edges.  Since the first $r$ rows of $M$ have zero interior
block, $R\mathbf1_b=0$.  Thus
$\operatorname{rank}R\le b-1<r$, and there is a nonzero
$u\in\ker R^{\mathsf T}$.  The vector $(u,0_b)$ is in the kernel of
$M$.  Its support lies wholly in permanent interior coordinates, which
future attachments and identifications never touch.  Extending this
vector by zeros gives a nonzero nullvector of every future grounded
matrix, since the grounded coordinate is retained rather than interior.
The determinant is therefore zero.
\end{proof}

The lemma permits an \emph{absorbing} zero response: all numerical
matrix entries may be discarded, although boundary geometry must still
be maintained.  The module records this zero separately from a missing
or interrupted computation.  This is a field-specific zero.  A zero
modulo one prime never proves that an integer determinant is zero.

The sufficient test $r\ge b$ is intentionally cheap.  A rank-deficient
$R$ can yield a permanent nullvector even when $r<b$; recognizing all
such cases would require another rank calculation.  The additional test
is unnecessary for the claimed size bound.  Every carried nonzero
response has
\[
                      r\le b-1,\qquad \dim M\le 2b-1.
\]

\subsection{Why every crossing needs only constantly many pivots}

Consider one admissible update.  Aggregate the old terminal rows and
columns under all identifications, insert the new crossing edge, and
move the newly closed black fragments into the interior.  Let their
number be $u$.  The old radical block remains zero, so the interior
matrix now has the form
\[
 A=\begin{pmatrix}0_r&C\\C^{\mathsf T}&D\end{pmatrix},
                     \qquad D\in\mathbb F^{u\times u}.
\]
The first $r$ rows are supported on $u$ columns; the final $u$ rows
contribute rank at most $u$.  Hence
\[
                         \operatorname{rank}A\le 2u.
\]
Symmetric nonsingular principal elimination removes exactly its pivot
rank and leaves an identically zero interior block.  It therefore
eliminates at most $2u$ coordinates, using at most $2u$ pivots.

Inserting one crossing creates at most four physical internal edges,
each supplying one black face connection.  Each newly closed black
fragment must contain a newly supplied connection, so $u\le4$.
Retaining the anchor can only reduce $u$.  Thus there are at most eight
one- or two-coordinate pivots per crossing, including exceptional
characteristics in which pivot profiles change.  No probabilistic
genericity or numerical stability assumption is used.

\subsection{All-stage complexity theorem}

\begin{theorem}[All-stage response compilation]
\label{thm:response-all-stage}
Let a spherical $n$-crossing source diagram have a prescribed crossing
order whose largest dart frontier is $W$.  Over any fixed field, all
nonempty suffix determinant responses can be compiled with
$O(n(W+1)^2)$ field operations.  Keeping all snapshots stores
$O(n(W+1)^2)$ field elements; streaming the snapshots uses
$O((W+1)^2)$ field elements in addition to $O(n)$ geometric storage.
An arbitrary checked boundary-matching query costs
$O((W+1)^3)$ field operations and $O(W\alpha(W)+W)$ combinatorial
operations.  These bounds include singular interiors.
\end{theorem}

\begin{proof}
For a nonzero response, $b\le W+1$ and $r\le b-1$.
At most two black vertices are introduced in the next update.
Identifications do not increase matrix dimension.  Matrix construction,
terminal aggregation, and each symmetric pivot therefore cost
$O((W+1)^2)$ field operations.  The preceding rank argument bounds
the number of pivots per crossing by an absolute constant.  A scan for
an available diagonal or off-diagonal pivot also examines at most
$O((W+1)^2)$ entries and is repeated only constantly often.  If the
absorbing-zero condition occurs, subsequent arithmetic is cheaper.
Summing over $n$ crossings proves the arithmetic claim.  Copying a
finished immutable response has the same quadratic bound.

Original edge pairing, source faces, checkerboard coloring, and the
insertion sequence are prepared in $O(n)$ graph work apart from
disjoint-set costs.  Union by size with path compression handles
$O(n)$ unions and $O(n+\sum_s w_s)$ finds.  Snapshot construction
visits only the current boundary, sorting its labels and roots once.
Its total combinatorial work is
\[
 O\!\left((n+\sum_s w_s)\alpha(n)
              +\sum_s w_s\log(w_s+2)\right).
\]
It never retraverses the whole suffix.  This separate statement avoids
silently treating union-find costs or index bit lengths as field
operations.  The geometric storage consists of the source darts,
two union-find structures, the active-crossing flags and the boundary
dictionary, all of size $O(n)$; a snapshot retains $O(W)$ labels.

For a query, boundary unions reconstruct exactly the completed open
faces and projection components.  The inherited coloring is checked,
and the sphere equation $F=n_s+2k$ is verified.  Closed fragments and
closed projection components contribute their stored counts.  There
are at most $W$ open black fragments.  Their quotient partition,
together with the fixed anchor when closed, grounds a matrix of
dimension at most $2b-2=O(W+1)$.  Dense field elimination costs the
stated cubic amount.
\end{proof}

Since $W\le 4n$, the construction also gives an unconditional
$O(n^3)$ field-operation bound for \emph{all-stage response compilation},
where independent dense cubic setup at every stage gives the coarser
$O(n^4)$ bound.  For $W=O(1)$ the new arithmetic bound is linear; for
$W=O(\log^c n)$ it is $O(n\log^{2c}n)$.  None of these statements
bounds the number of observations or the work of the relative-complex
scanner that requests them.

For a variable prime $p$, multiplication costs and inversions must be
restored to the field-operation bound.  With $\ell=\lceil\log_2p\rceil$,
the construction uses $O(n(W+1)^2)$ operations on $O(\ell)$-bit field
elements and only $O(n)$ inversions.  The implementation validates
primes below $2^{31}$, including two, so its individual field elements
have a fixed machine-word bound.  More primes multiply the compilation
cost; the theorem does not claim constant-cost reconstruction of an
arbitrarily large integer cofactor.

\subsection{Implementation and integration boundary}

The new file \path{fastunknot/streaming_response.py} supplies immutable
\texttt{StreamingResponse} and \texttt{SuffixResponse} objects, an
iterator over suffixes in reverse stage order, and a convenience
function returning every snapshot indexed by forward prefix stage.
The generic arithmetic API rejects new edges whose endpoints are no
longer retained terminals.  The geometric compiler additionally checks
every black-fragment merge against its live terminal set, so an
accidental identification of an eliminated face is an explicit error.

Snapshot completion queries check exact boundary coverage, inherited
color compatibility, and spherical genus.  A disconnected projection
has raw reduced Jones value zero at $i$.  A connected query uses the
phase $(-i)^{B+v-1}$, where $B=\sum_e b_e$: the crossing bit is
$b_e=1$ for black sectors $[0,2]$ and $b_e=0$ for black sectors $[1,3]$,
and the edge weight is $(-1)^{b_e}$. This sum retains contributions of
Laplacian loops. The number $v$ is the actual completed black-face count.
The anchor is included in this count exactly once.  Modulo a prime the
result is stored as a formal Gaussian pair; it does not require
choosing a square root of $-1$ in the field.

All response updates and queries poll an optional cancellation/work
callback.  An exception propagates before a partially completed
snapshot is published.  Existing snapshots cannot be modified by an
update.  The constructor's source validation is also interruptible.
There is no new default setting, fallback-policy change or positive
unknot claim.

The iterator accepts an optional frontier allowance and stops before a
response at an over-wide stage is built.  The all-snapshot convenience
function also accepts a retained-field-element allowance, charging the
sum of squared stored response dimensions before retaining each completed
snapshot.  Exhaustion raises a distinct \texttt{ResponseBudget} exception;
the convenience function returns no partial compiled tuple.  The storage
allowance covers retained field entries, not source metadata or a current
update's temporary matrix.  It is not a hard process-memory limit.

The primary integration opportunity is an observer that expects to
query many stages of a fixed order.  It can prepare responses once
and query the corresponding stage without repeating interior
elimination.  Conversely, a scan that is likely to terminate after
one inexpensive early observation should not automatically prepare
every suffix.  An adaptive dispatcher must account for the all-stage
setup cost.  No such dispatcher or end-to-end speedup is claimed here.

\subsection{Independent checks and measurements}

The reproducible integer-cofactor audit accepts 53 knot braid diagrams
from 240 seeded attempts.  It considers every one of their 610 suffixes
under a recorded shuffled order, over the primes $2,3,5,101$.
There are 2440 field responses and 19520 independent quotient-cofactor
comparisons, all agreeing.  The sample includes 493 singular responses,
19 absorbing-zero responses, and 72 streams that perform a
two-coordinate pivot.  The largest retained radical has dimension four.
The reference reconstructs complete cut-face graphs by graph components
and computes integer Bareiss cofactors; it does not reuse the incremental
fragment union-find or the field-response matrices.

Eight focused tests additionally exercise actual scanner matchings and
their raw Gaussian phases, mirrors, randomly generated graph insertion
streams with identifications and forgetting, the characteristic-two
singular path example, protected zero responses, immutable updates,
source and mid-update cancellation, and the explicit frontier/storage
allowances.  No first-jet or recognition driver is imported by the new
response compiler.

The timing family is the natural crossing order of the positive
three-strand braid $(\sigma_1\sigma_2)^{n/2}$ for the listed $n$.
The closures are knots since $n/2$ is coprime to three.  Both arms
prepare every suffix of the same diagram.  The dense control independently
rebuilds each cut-face graph and uses the same symmetric pivot arithmetic;
the streaming arm retains elimination between suffixes.  Original
coloring is computed once per arm.  Source preparation, geometry and
all response setups are included.  One excluded warmup checks five
arbitrary quotient queries per suffix.  Shuffled A/A/B rounds include
two identical dense controls.  The first four rows use seven rounds;
the last uses three.

\begin{center}
\begin{tabular}{r|r|r|r|r}
$n$ & Dense setup (ms) & Streaming (ms) & Ratio & Dense A/A\\\hline
16 & 1.308 & 0.522 & 2.51 & 1.076\\
32 & 5.584 & 1.123 & 4.97 & 1.155\\
64 & 27.180 & 2.193 & 12.39 & 0.993\\
128 & 267.265 & 5.241 & 50.99 & 0.992\\
256 & 3615.101 & 14.061 & 257.10 & 0.906
\end{tabular}
\end{center}

The large ratios reflect substantial interior reuse.  At $n=256$,
253 of 256 stages have a nontrivial interior and 84 have a nonzero
retained radical.  The dense common graph reaches 129 black vertices;
the largest matrix constructed by streaming has dimension five.
Each dense arm performs 23789518 symmetric elimination updates across
all suffixes, compared with 928 for streaming. At all natural-order sizes
the largest dart frontier is six.  Compiling all 8192 suffixes of the
corresponding larger braid takes 0.4155 seconds in one additional
streaming-only run, performs 30024 elimination updates, and stores
92825 field elements.  That run has no dense timed comparator.

Recorded shuffled orders at 16, 32 and 64 crossings give smaller ratios
of 1.69, 2.50 and 3.73.  This sensitivity is expected: a poor order
increases the actual retained boundary.  Small-instance timings and
the 256-crossing dense controls exhibit visible noise; the table
reports their controls rather than treating milliseconds as exact.
The operation-count theorem is independent of these wall-clock ratios.

These knots are readily decided by existing recognizer filters.
The measurements establish a response-compilation improvement, not a
whole-recognizer speedup or a difficult family for unknot recognition.
Their significance relative to earlier boundary benchmarks is that
substantial, sometimes singular interiors are eliminated once across
genuinely changing suffixes.  All raw input words, orders, timings,
counts and source hashes are retained in
\path{fast/results/streaming_response*.json}; the source hashes in
all three final JSON files match the delivered response module.

\subsection{Further questions specific to this construction}

\begin{enumerate}
\item Can a cost-based observer choose a prefix of the reverse compilation
      or a collection of checkpoints that gives a competitive bound relative
      to the better of independent setup and full precompilation?
\item Can the same face filtration compile the two Euler boundary path
      summaries and determinant response in one pass, so stage geometry and
      cancellation checks are shared?
\item Can a separator tree replace a linear crossing order while preserving
      singular residual information, giving local updates after a small
      diagram change?  The present no-future-interior-identification contract
      must be replaced by a proved tree-boundary contract.
\item Which real observer traces contain enough nontrivial interior work
      at enough stages for this compilation theorem to improve complete
      recognition?  Easy torus-braid families isolate the arithmetic gain
      but cannot answer this question.
\item Can boundary response snapshots be shared across several primes or
      several crossing orders without assuming a common pivot profile?
      Exceptional primes require exact handling of changing radicals.
\end{enumerate}
